% !TeX root = ../main.tex
\clearpage
\section{Discussion}

This thesis draws on a multi-year record of electric organ discharges (EODs) from captive \textit{Electrophorus voltai}.
The analysis covers pulse shapes, pulse rates across several timescales, coarse space use along an electrode line, and the response of pulse rate to feeding. 
The dataset contains more than thirteen million accepted EODs and spans approximately 14 TiB of raw recordings. 
A big share of this work also includes the hardware development and field work that made these recordings possible. 
In that sense, the thesis is as much a methodological contribution as an ethological one.
The main ethological value of the dataset is that its duration makes it possible to identify changes in EOD activity that would be difficult to distinguish from short-term variation in a conventional behavioral observation.

Across the Berlin dataset, the classifier assigns accepted EODs to normal, wide, and double labels. 
These labels are useful working categories for summarising waveform variation, but neither the half-width distributions nor the Robust PCA support three cleanly separated morphological islands. 
The overall picture is closer to a continuum of waveform morphology, across which pulses can vary in structural properties such as duration, shoulder strength, and relative peak timing. 
That continuum makes hard boundaries between classes difficult to justify and means that any discrete label scheme remains an approximation. 
At the same time, the extremes of this continuum remain striking: short symmetric single peaks, elongated single peaks with a shoulder, and clear two-peak envelopes differ enough that describing them is a necessary first step toward asking why EOD shape is so diverse and how it is manipulated. 
Candidate explanations include voluntary or involuntary behavioural state as well as external factors such as stress or physical conditions; the present data do not yet distinguish among them.

Previous work on \textit{Electrophorus} has mainly distinguished discharge forms by amplitude and behavioural context. 
For example low-amplitude search pulses, high-amplitude hunting volleys \parencite{catania2019astonishing}, and a third form that differs primarily in amplitude rather than in waveform shape \parencite{xu2021third}.
The present analysis keeps low- and high-voltage discharges in one accepted set and then groups them by waveform morphology, independently of amplitude or of whether they occur in search activity or in high-rate volleys.
Rate, feeding, and position analyses use that same inclusive set.
In that sense, the Berlin dataset provides, to our knowledge, the first large-scale, long-term characterization of this shape variation in \textit{Electrophorus}.
The labels describe waveform shape and do not, by themselves, identify behavioural acts. 
This distinction is important throughout the interpretation of the results.

The higher PC views, in which doubles split into two clumps, and the two-template fits of the median double, together suggest that the double label itself is not a single stereotyped waveform.
Recovering the trough requires components broader than the native normal prototype, so the double is consistent with a short doublet of broadened normal-like components rather than an unmodified normal repeated about a millisecond later.
That morphology should also be distinguished from the high-voltage doublets described during prey capture \parencite{catania2014shocking}: the Berlin doubles occur across a range of signal strengths, and they are scored as one two-peaked event rather than as two separate pulses.
One speculative physiological route would involve staggered contributions from different electric organs under a single pacemaker command; ion-channel work in weakly electric fish shows that spike duration can shape EOD width \textcite{mcanelly2000coregulation}, and the membrane physiology of the organ of Sachs remains relevant background \parencite{keynes1953membrane}, but neither supplies a direct mechanism here.
Possible behavioural associations of these morphological extremes are taken up below in the feeding analysis and in the late-2025 activity pattern.

\subsection{Feeding-related activity}

The feeding analysis provides the clearest behavioural association in the dataset. 
Overall pulse rate rises sharply in the minutes around logged feeding, and the peri-feeding trajectory begins to climb before the logged feeding time.
This pattern also helps explain the morning maximum in the 24~h activity plot. 
The timing is consistent with the animals increasing their pulse rate before food arrives, which could reflect anticipation of the feeding schedule or responses to cues associated with the keeper and feeding routine.

Wide pulses follow the same peri-feeding shape at lower absolute rates.
This association suggests that wide pulses are involved in the elevated EOD activity surrounding feeding. 
It is tempting to relate this finding to hunting, since electric eels use their electric system extensively during prey detection and capture \parencite{catania2014shocking, catania2015highvoltage, catania2015electric, catania2019astonishing}. 
However, the present recordings were not linked to individual prey encounters or predatory strikes. 
The feeding response therefore provides evidence for an association between wide pulses and feeding-related activity, but it does not establish that wide pulses are hunting signals.
That limitation is unsurprising here: the eels were fed dead fish or fish pieces rather than live prey, so a classic live-prey hunting signature was not expected.

Double-pulse means stay close to baseline in the feeding windows, yet the minute-level distributions still differ and a brief peri-event rise is visible.
Rather than isolating double pulses as a feeding- or hunting-specific signal, this pattern is more consistent with waveforms that change only slowly while the underlying electrolocation function of the EOD is retained across shape classes. 
The morphological similarity to hunting doublets described during prey capture \parencite{catania2014shocking} is therefore interesting, but similarity alone is insufficient to assign the Berlin double-pulse class to hunting behaviour. 
Directly video-linked recordings of feeding events would be a useful next step for resolving this question.

\subsection{Temporal patterns and the late-2025 increase}

The 24~h activity profile contains a pronounced morning increase, which coincides closely with the regular feeding schedule. 
On the calendar-month scale, overall and wide-pulse rates rise in the later part of the year, but the chronological view shows that this is driven mainly by a concentrated late-2025 episode rather than by a pattern that recurs across years.
Wide pulses are almost absent earlier in the record and reach their clearest increase in that window; doubles remain rarer but also rise toward the end of the recordings.
A simple seasonal explanation is therefore unlikely: if the increase were seasonal, a similar rise would be expected in corresponding months of more than one year.
This does not rule out environmental or husbandry cues entirely.
During the same broader late-2025 window, logged tank conductivity sat toward the lower end of the measured range, but November and December 2025 themselves have little overlapping conductivity coverage, so that coincidence is not a demonstrated driver.
The present data likewise do not support reading the November 2025 increase as a recurring rainy-season pattern.

The late-2025 increase becomes more interesting when considered together with space use and anecdotal behaviour.
Session-wise position estimates move toward the dark boundary in the same period, while overall occupancy already favours the tank ends and the darker mid-line electrodes more than the bright mid-line.
The dark section contains root structures that could support breeding-related behaviour, and reproductive behaviour in the wild is closely associated with structured habitats and root masses \parencite{assunccao1995reproduction, bastos2020historia}.
The spatial data cannot show that the animals used that region for reproductive purposes; they only show that space use changed alongside the activity increase.
Caretaker notes from the same period describe the eels remaining closer to one another than usual and actively seeking proximity, particularly unusual for the female.
Those notes were not scored with a standardized protocol, so they are circumstantial context rather than independent quantitative evidence.

Taken together, the special-pulse increase, the shift toward the darker region, and the proximity observations are consistent with a change in social state.
A reproductive reading is especially interesting given how little is known about electric eel courtship, but the recordings do not identify mating, courtship, nest building, spawning, or any other specific reproductive act, and the position estimate cannot separate individuals when several eels contribute to a recording.
The late-2025 pattern therefore supports a testable association between special-pulse activity and an unusual social or potentially reproductive state; it does not show that the special pulses are reproductive signals.
The subsequent deaths of the female in late December and of the male about four months later fall in the same broader unusual period and, if a reproductive condition such as egg retention were confirmed, would add circumstantial context---but cause of death cannot be established from the present dataset and must stay separate from the quantitative evidence.

\subsection{Environmental variables and long-term changes}

Environmental variables provide a less clear alternative explanation for the observed activity patterns. 
Temperature and conductivity showed no robust relationships with overall, wide, or double-pulse rates, and the one significant Pearson association (double-pulse rate vs temperature) failed under Spearman ranking.
The relatively narrow temperature range in the tank also limits the ability of this analysis to detect environmental effects. 
A lack of a correlation within the measured range should not be interpreted as evidence that EOD activity is generally independent of temperature. 
Temperature-related changes of EOD waveform have been reported in weakly electric fish, including \textit{Brachyhypopomus} and \textit{Gymnotus} \parencite{silva1999water, ardanaz2001temperature}.
The temporal coincidence noted above between the late-2025 special-pulse increase and comparatively low tank conductivity therefore remains only a qualitative observation: it does not override the absence of a robust rate--conductivity correlation across the logged months.

The long-term decline in median all-pulse rate across yearly bins is at least partly explained by how the recordings were sampled and by hardware limits.
Early 2024 highs are dominated by short morning feeding sessions with high rates, whereas from July onward longer overnight runs contribute more to the monthly medians and pull them down.
Session notes from mid-July 2024 also report that the electrode line was dead from channel 9 onward, leaving only 12 channels, which could undercount pulses whenever the animals were toward the dark end of the line; processed recordings show all 16 channels again from 7~August~2024 onward.
Aging, habituation, and changes in the number of animals remain possible additional factors.
Dual-line pulse counts are already weighted from 25~November~2025, so the yearly decrease is not simply an effect of adding a second electrode line. 
The partial coverage of 2023 and 2026 also limits direct comparison between those years.
Importantly, the special pulse classes do not simply follow the same long-term decline as all pulses, which suggests that the reduction is not a uniform quieting of every discharge class.

\subsection{Spatial behaviour and limitations of the position estimate}

Spatial exploration of the Berlin dataset is necessarily coarse. 
Head position is estimated from the electrode with the largest positive peak.
This estimate is most reliable when one animal dominates a short recording segment and is oriented reasonably along the electrode line. 
It becomes less reliable when several animals discharge simultaneously or when an animal lies parallel to the array.
Morning movement toward the bright end coincides with the feeding response, and that end is also the keeper and feeding side of the tank.
The spatial scale of the tank also limits comparison with the natural environment. 
Improvements to the localisation procedure, including further development of the eeltracker processing chain, would make it possible to reconstruct trajectories rather than only occupancy frequencies. 
The FLONA data were mainly collected with electrode grids rather than a line, which opens the possibility of a more complete spatial reconstruction in the natural environment.

\subsection{Captivity and comparison with the natural environment}

The distinction between the captive and natural environments is important for interpreting the absence of strong periodic behaviour as well as the late-2025 increase. 
Wild electric eels are exposed to environmental cues that are difficult to reproduce in captivity, including changes in hydrology, habitat structure, and food availability \parencite{assunccao1995reproduction, bastos2020historia}.

Although the husbandry conditions were adjusted to provide potential cues for breeding, the captive animals were not exposed to the full set of environmental changes experienced in their natural habitat. 
A lack of a strong periodic signal in the Berlin recordings should therefore not be interpreted as evidence that electric eel behaviour itself lacks periodicity.

The Berlin dataset is particularly valuable because its controlled setting and long duration make it possible to identify waveform classes, obtain enough rare double pulses to examine their morphology, and compare pulse rates around repeated feeding events. 
At the same time, the captive setting cannot resolve behaviours that are unlikely to occur in the tank, including migration, natural nest building, natural prey encounters, and the group hunting behaviour reported in the field \parencite{bastos2021social}. 
The Berlin dataset therefore provides a controlled baseline rather than a substitute for observations in the natural habitat.

The field recordings collected during the September 2025 expedition provide the most direct next step. 
Applying the same waveform classification and activity analysis to the FLONA and Ducke recordings would allow the special pulse classes identified in Berlin to be examined under natural conditions.
This comparison could determine whether their relative rates, pulse-rate timescales, feeding-related activity, and periodic patterns are similar between captivity and the wild. 
It could also reveal patterns that are absent in the captive dataset because the corresponding environmental or social conditions were not present.

This comparison is important for interpreting the behavioural significance of the present findings. 
If similar waveform and activity patterns occur in the wild, this would strengthen the argument that they represent biologically meaningful aspects of \textit{E. voltai} behaviour rather than responses specific to captivity. 
Conversely, differences between the datasets would help identify which aspects of the observed activity are associated with the captive environment.

The recording system itself is an important part of what makes these analyses possible. 
Pulse amplitudes measured on the grid fall within the range of tens of volts reported previously by \textcite{xu2021third}, providing a useful check that passive non-invasive electrodes recover the expected voltage scale. 
The TeeGrid, TeeRec, and Teensy amplifier system, together with the waterproofing and battery solutions developed for the field work, provide a basis for extending the same recording approach beyond the captive setting.
For software availability and hardware details, see Section~\ref{sec:recording_hardware}.

\subsection{Conclusions and Outlook}

The Berlin recordings provide a long-term captive baseline in which normal, wide, and double labels mark reproducible extremes of a morphological continuum, feeding elevates overall and wide-pulse activity with a pre-feeding rise, and coarse localisation shows non-uniform tank use.
The most distinctive pattern is the late-2025 rise in special-pulse activity together with a shift toward the dark boundary and anecdotal proximity between the animals: consistent with a change in social state, and a reproductive hypothesis worth testing, but not evidence that the special pulses are reproductive signals.

Several questions remain unresolved. 
The waveform labels should not yet be interpreted as behavioural categories, and the double pulses in particular require recordings linked to individual animals, amplitude, and behaviour before they can be related to hunting or other specific actions. 
A related next experiment would be to relate special-pulse labels to high- versus low-voltage discharges explicitly: because the present classification ignores amplitude, it remains open whether wide or double shapes are enriched in one voltage regime, or whether their occurrence is mainly context-dependent rather than voltage-dependent. 
The cause of the long-term decline in overall pulse rate also remains unclear, as does the extent to which the narrow range of environmental variables in the Berlin tank constrains the observed patterns.

The most important next step is to apply the same classification and activity pipeline to the FLONA and Ducke recordings from the natural environment.
Captivity supplied the duration and controlled context needed to discover these patterns; the field data are needed to determine their ecological meaning.
The main contribution of this thesis is therefore not that it resolves the social or reproductive behaviour of electric eels, but that long-term passive EOD recordings can reveal waveform variation and activity changes that are difficult to detect through visual observation alone.
